% A self-contained section for the companion research note.
% Requires amsmath, amssymb, amsthm, booktabs, hyperref and theorem environments.
% Bibliography keys are supplied in network_references.bib.
\section{Optimal networks for graph-coupled zeros}
\label{net:sec:optimal}

The article ends its section on zeros with an explicit optimization question:
given the degrees of the input polynomials and a fixed total conductance,
which graph minimizes the largest possible displacement of the coupled
zeros? We solve that question when all pairs of vertices are permitted to
be joined. We also solve the corresponding mean-displacement problem.
The two optimal networks usually differ. Both are determined by the input
degrees alone, independently of every input root.

The optimization arguments use elementary matrix inequalities. The link to
zeros is the article's compression and electrical-transport identity
\cite[\texttt{zer:thm:electrical}]{ProveItArticle}. Conductance optimization
and fractional graph Laplacians have an established literature
\cite{GhoshBoydSaberi2008,Schweitzer2022}; the results below resolve the
particular coupling problem posed in the article, with explicit edges,
equality conditions and stability estimates. They are unconditional
statements about real-rooted polynomials. Dirichlet GRH is not needed for
these finite polynomial assertions.

\subsection{The degree matrix and the two objectives}

Fix $n\ge2$, positive integers $d_1,\ldots,d_n$, and a budget $W>0$.
The numbers $d_i$ are the degrees of the input polynomials, not the graph
degrees. Write
\begin{equation}\label{net:eq:degree-matrix}
 d=(d_1,\ldots,d_n)^{\mathsf T},\qquad N=\sum_i d_i,\qquad
 C=\operatorname{diag}(d)-\frac{dd^{\mathsf T}}{N},\qquad
 H=\mathbf1_n^{\perp}.
\end{equation}
Set
\begin{equation}\label{net:eq:TS}
 T=\operatorname{tr}C=N-\frac{\sum_i d_i^2}{N},\qquad
 S=\operatorname{tr}\sqrt C,\qquad b=2W.
\end{equation}
The square root is the positive semidefinite matrix square root. In
particular, it vanishes on $\mathbf1_n$.

An admissible graph has nonnegative conductances $w_{ij}=w_{ji}$,
$w_{ii}=0$, a connected support, and $\sum_{i<j}w_{ij}=W$. Its Laplacian
$L$ satisfies $L_{ij}=-w_{ij}$ for $i\ne j$, $L\mathbf1_n=0$, and
$\operatorname{tr}L=b$. Thus $L_H=L|_H$ is positive definite. Missing
edges are allowed; the underlying topology is part of the optimization.
Parallel edges are combined by adding their conductances.

The matrix $C$ itself is the Laplacian of the complete graph with edges
$d_i d_j/N$. Indeed,
\begin{equation}\label{net:eq:C-energy}
 x^{\mathsf T}Cx=\frac1N\sum_{i<j}d_i d_j(x_i-x_j)^2.
\end{equation}
Consequently $C_H$ is positive definite and $T,S>0$.

For monic real-rooted polynomials $p_i$ of degrees $d_i$, let
\begin{equation}\label{net:eq:forest}
 \mathcal P_L(z)=\sum_F
 \left(\prod_{e\in F}w_e\right)
 \left(\prod_{i\in\mathcal R(F)}p_i'(z)\right)
 \left(\prod_{i\notin\mathcal R(F)}p_i(z)\right),
\end{equation}
where the sum is over rooted spanning forests and $\mathcal R(F)$ contains
one chosen root of each component. Repeated roots and roots common to
different polynomials are permitted. Let $\lambda_1\le\cdots\le
\lambda_{N-1}$ be the zeros of $\mathcal P_L$ and let
$\beta_1\le\cdots\le\beta_{N-1}$ be the zeros of
$(\prod_i p_i)'$.

\begin{lemma}[The graph-dependent objectives]\label{net:lem:objectives}
All the $\lambda_k$ are real, and $\lambda_k\le\beta_k$. Define
\begin{align}
 E_{\infty}(L)
   &=\lambda_{\max}\!\left(C_H^{1/2}L_H^{-1}C_H^{1/2}\right),
       \label{net:eq:infty}\\
 E_1(L)&=\operatorname{tr}(L^+C), \label{net:eq:one}
\end{align}
where $L^+$ is the Moore--Penrose inverse. Then
\begin{equation}\label{net:eq:zero-objectives}
 \max_k(\beta_k-\lambda_k)\le E_{\infty}(L),\qquad
 \sum_{k=1}^{N-1}(\beta_k-\lambda_k)=E_1(L).
\end{equation}
The first bound is attained when $p_i(z)=(z-c)^{d_i}$ for any real $c$.
Thus $E_{\infty}(L)$ is the worst possible largest displacement over all
permitted inputs, while $E_1(L)/(N-1)$ is the exact mean displacement for
every input.
\end{lemma}

\begin{proof}
List the input roots, with multiplicity, in vertex blocks. Let
$\Lambda$ be their $N\times N$ diagonal matrix, let $V_{ai}$ be one when
root index $a$ belongs to block $i$ and zero otherwise, and put
$Q=I_N-\mathbf1_N\mathbf1_N^{\mathsf T}/N$.
On $\mathcal V=\mathbf1_N^\perp$ define
\[
 A=(Q\Lambda Q)|_{\mathcal V},\qquad
 B=(QVL^+V^{\mathsf T}Q)|_{\mathcal V}.
\]
The compression identity in the article is
\[
 \mathcal P_L(z)=N\tau(L)\det_{\mathcal V}(zI-A+B),\qquad
 \frac1N\left(\prod_i p_i(z)\right)'
       =\det_{\mathcal V}(zI-A),
\]
with $\tau(L)>0$ the spanning-tree sum. For completeness, the first
identity follows by substituting $x_i=p_i'/p_i$ into the matrix-forest
identity $\det(L+\operatorname{diag}x)=\sum_F w(F)\prod_{i\in
\mathcal R(F)}x_i$, multiplying by $\prod_i p_i$, and taking the Schur
complement of the positive part of $L$ in
\[
 \begin{pmatrix}zI_N-\Lambda&V\\-V^{\mathsf T}&L\end{pmatrix}.
\]
The remaining null direction is $\mathbf1_n/\sqrt n$, whose image under
$V$ has squared norm $N/n$; $\det{}'L=n\tau(L)$ gives the factor
$N\tau(L)$. The derivative identity follows from the cofactor formula for
a compression to $\mathbf1_N^\perp$.

Since $B\succeq0$, the min--max principle gives the ordering of the zeros
and bounds every displacement by $\|B\|_{\mathrm{op}}$. Their sum is
$\operatorname{tr}B$. Now
\[
 V^{\mathsf T}QV=C.
\]
The nonzero eigenvalues of $B$ equal those of
$(L^+)^{1/2}C(L^+)^{1/2}$, by applying the equality of the nonzero spectra
of $XX^{\mathsf T}$ and $X^{\mathsf T}X$ to
$X=QV(L^+)^{1/2}$. On $H$, the latter has the same eigenvalues as
$C_H^{1/2}L_H^{-1}C_H^{1/2}$. This proves
\eqref{net:eq:zero-objectives}. If all roots are $c$, then $A=cI$, and
the displacements are exactly the eigenvalues of $B$.
\end{proof}

The trace objective also has the electrical form
\begin{equation}\label{net:eq:resistances}
 E_1(L)=\frac1N\sum_{i<j}d_i d_jR_{ij},\qquad
 R_{ij}=(e_i-e_j)^{\mathsf T}L^+(e_i-e_j).
\end{equation}
It follows by expanding the right side and using \eqref{net:eq:degree-matrix}.

\subsection{The sharp optima}

\begin{theorem}[Optimal worst displacement and optimal mean displacement]
\label{net:thm:optima}
Among all admissible graphs,
\begin{align}
 \min_L E_{\infty}(L)&=\frac{T}{2W},&
 L_{\infty}&=\frac{2W}{T}C, \label{net:eq:optimal-infty}\\
 \min_L E_1(L)&=\frac{S^2}{2W},&
 L_1&=\frac{2W}{S}\sqrt C. \label{net:eq:optimal-one}
\end{align}
Each minimizing Laplacian is unique. Both are Laplacians of complete graphs
with strictly positive conductances. The first has the explicit edges
\begin{equation}\label{net:eq:product-edges}
 w^{(\infty)}_{ij}
   =\frac{W d_i d_j}{\sum_{a<b}d_a d_b}\qquad(i\ne j).
\end{equation}
In particular, the smallest possible worst displacement is
\[
 \frac{\sum_{i<j}d_i d_j}{WN},
\]
and the smallest possible mean displacement is $S^2/[2W(N-1)]$.
The same optima hold with the budget constraint $\sum_{i<j}w_{ij}\le W$.
\end{theorem}

\begin{proof}
All inverse matrices in this proof act on $H$. The condition
$E_{\infty}(L)\le t$ is equivalent to
\begin{equation}\label{net:eq:order}
 C\preceq tL.
\end{equation}
To see this, $C_H^{1/2}L_H^{-1}C_H^{1/2}\preceq tI$ is equivalent to
$L_H^{-1}\preceq tC_H^{-1}$; inversion of positive definite matrices
reverses the order. Taking the trace of \eqref{net:eq:order} gives
$T\le bt$. The matrix $L_{\infty}=bC/T$ is admissible, and realizes
equality. If $t=T/b$, then $tL-C$ is positive semidefinite with trace zero,
so $tL=C$. This proves the first optimum and its uniqueness.

For the second, the Hilbert--Schmidt Cauchy--Schwarz inequality gives
\begin{align*}
 S^2
 &=\left\langle L_H^{1/2},L_H^{-1/2}C_H^{1/2}\right\rangle_F^2\\
 &\le \operatorname{tr}(L_H)\,
       \operatorname{tr}(C_H^{1/2}L_H^{-1}C_H^{1/2})
   =bE_1(L).
\end{align*}
Equality holds precisely when
$L_H^{-1/2}C_H^{1/2}=(S/b)L_H^{1/2}$, or
$L_H=(b/S)C_H^{1/2}$. Lemma~\ref{net:lem:edges} below proves that this
matrix is a complete-graph Laplacian. It has trace $b$, so it is
admissible and gives the unique minimum. Formula
\eqref{net:eq:product-edges} follows from
$T=2\sum_{i<j}d_i d_j/N$.

Finally, scaling all conductances by $a>0$ divides both objectives by $a$.
Any graph using strictly less than the available budget can therefore be
improved by scaling up.
\end{proof}

The minimax graph equalizes the nonzero displacements in the equality
example: when $p_i(z)=(z-c)^{d_i}$, its coupled zeros are $c-T/(2W)$
with multiplicity $n-1$ and $c$ with multiplicity $N-n$. The mean-optimal
graph generally distributes these displacements unequally, reducing their
sum.

\subsection{Explicit edges of the square-root network}

\begin{lemma}[Positive integral and contrast bound]\label{net:lem:edges}
For $t>0$ put $h(t)=\sum_{k=1}^n d_k/(t+d_k)$. Then, for $i\ne j$,
\begin{equation}\label{net:eq:integral-edges}
 w^{(1)}_{ij}
 =\frac{2W d_i d_j}{\pi S}
   \int_0^{\infty}
   \frac{dt}{\sqrt t\,(t+d_i)(t+d_j)h(t)}.
\end{equation}
The integral converges and is strictly positive. If $i,j,k$ are distinct
and $d_j>d_k$, then
\begin{equation}\label{net:eq:contrast}
 1<\frac{w^{(1)}_{ij}}{w^{(1)}_{ik}}<\frac{d_j}{d_k},
 \qquad
 \frac{w^{(\infty)}_{ij}}{w^{(\infty)}_{ik}}=\frac{d_j}{d_k}.
\end{equation}
If $d_j=d_k$, both ratios are one. Thus the mean-optimal network retains
the ordering of the degree weights but strictly reduces each nontrivial
contrast within a row of edges.
\end{lemma}

\begin{proof}
The scalar identity
\[
 \sqrt x=\frac1\pi\int_0^\infty
 t^{-1/2}\frac{x}{t+x}\,dt\qquad(x\ge0)
\]
follows for $x>0$ by substituting $t=xu^2$; it is immediate for $x=0$.
The spectral theorem therefore gives the entrywise convergent identity
\begin{equation}\label{net:eq:sqrt-resolvent}
 \sqrt C=\frac1\pi\int_0^\infty t^{-1/2}C(tI+C)^{-1}\,dt.
\end{equation}
Write $A_t=tI+\operatorname{diag}(d)$. The rank-one inverse formula yields
\[
 (tI+C)^{-1}
 =A_t^{-1}
 +\frac{A_t^{-1}dd^{\mathsf T}A_t^{-1}}
        {N-d^{\mathsf T}A_t^{-1}d},\qquad
 N-\sum_k\frac{d_k^2}{t+d_k}=t h(t).
\]
Consequently, for $i\ne j$,
\[
 ((tI+C)^{-1})_{ij}
 =\frac{d_i d_j}{t(t+d_i)(t+d_j)h(t)}.
\]
Taking off-diagonal entries in
$C(tI+C)^{-1}=I-t(tI+C)^{-1}$ gives
\eqref{net:eq:integral-edges}. Near zero the integrand is
$O(t^{-1/2})$; near infinity it is $O(t^{-3/2})$, since
$h(t)\sim N/t$. It is positive everywhere. Moreover
$\sqrt C\mathbf1_n=0$, so these negative off-diagonal entries and their
row sums make $\sqrt C$ a complete-graph Laplacian. Its scale in
\eqref{net:eq:optimal-one} has total conductance $W$.

For the contrast assertion, put
$f_j(t)=d_j/(t+d_j)$. Cancelling the common factors in the integrals
shows that $w^{(1)}_{ij}/w^{(1)}_{ik}$ is the average of
\[
 \frac{f_j(t)}{f_k(t)}
 =\frac{d_j}{d_k}\frac{t+d_k}{t+d_j}
\]
against the positive finite measure with density
$t^{-1/2}f_i(t)f_k(t)/h(t)$, normalized to have total mass one.
For every $t>0$ this ratio is strictly between $1$ and $d_j/d_k$ when
$d_j>d_k$. Integration preserves the strict inequalities.
\end{proof}

\begin{corollary}[When the two optima agree]\label{net:cor:agreement}
For $n\ge3$, $L_1=L_{\infty}$ if and only if
$d_1=\cdots=d_n$. For $n=2$, they always agree: the only admissible graph
is the single edge of conductance $W$.
\end{corollary}

\begin{proof}
Equality of the two normalized matrices means $\sqrt C$ is proportional
to $C$. Hence all $n-1$ positive eigenvalues of $C$ coincide, and
$C=\gamma(I-\mathbf1_n\mathbf1_n^{\mathsf T}/n)$ for some $\gamma>0$.
Comparing off-diagonal entries gives $d_i d_j=\gamma N/n$ for every
$i\ne j$. When $n\ge3$, compare the pairs $(i,j)$ and $(i,k)$ to obtain
$d_j=d_k$, and repeat. Conversely equal $d_i$ give exactly such a scalar
multiple of the projection. The two-vertex assertion is immediate.
\end{proof}

\paragraph{An exact unequal-degree example.}
Take $(d_1,d_2,d_3)=(1,1,4)$ and $W=1$. The positive eigenvalues of $C$
are $1$ and $2$, so $T=3$ and $S=1+\sqrt2$. The complete comparison is
\begin{center}
\begin{tabular}{@{}lcc@{}}
\toprule
Quantity&Minimax network $L_{\infty}$&Mean network $L_1$\\
\midrule
$w_{12}$&$1/9$&$\dfrac{3-\sqrt2}{3(1+\sqrt2)}$\\[8pt]
$w_{13}=w_{23}$&$4/9$&$\dfrac{2\sqrt2}{3(1+\sqrt2)}$\\[8pt]
$E_{\infty}$&$3/2$&$(2+\sqrt2)/2$\\[4pt]
$E_1$&$3$&$(3+2\sqrt2)/2$\\
\bottomrule
\end{tabular}
\end{center}
For the mean displacement, divide the last row by $N-1=5$.
Thus minimizing the largest displacement and minimizing its mean are
different decisions even for three vertices.

\subsection{Stability and scope}

Write $\|A\|_{S_1}$ for the nuclear norm, the sum of the singular values,
and $\|A\|_F$ for the Frobenius norm. The following estimates quantify
the uniqueness in Theorem~\ref{net:thm:optima}.

\begin{theorem}[Stability of both designs]\label{net:thm:stability}
For every admissible $L$ and every $\varepsilon\ge0$:
\begin{enumerate}
\item If $E_{\infty}(L)\le(1+\varepsilon)T/b$, then
\begin{equation}\label{net:eq:stability-infty}
 \|L-L_{\infty}\|_{S_1}\le\frac{2b\varepsilon}{1+\varepsilon},\qquad
 \|L-L_{\infty}\|_{\mathrm{op}}\le\frac{b\varepsilon}{1+\varepsilon}.
\end{equation}
\item There is an exact identity
\begin{equation}\label{net:eq:residual}
 E_1(L)-\frac{S^2}{b}
 =\left(\frac{S}{b}\right)^2
 \left\|L_H^{-1/2}(L_H-(L_1)_H)\right\|_F^2.
\end{equation}
In particular, if $E_1(L)\le(1+\varepsilon)S^2/b$, then
\begin{equation}\label{net:eq:stability-one}
 \|L-L_1\|_{S_1}\le b\sqrt\varepsilon.
\end{equation}
\end{enumerate}
\end{theorem}

\begin{proof}
Let $t=E_{\infty}(L)$, $t_*=T/b$, and
$\eta=t/t_*-1\in[0,\varepsilon]$. By \eqref{net:eq:order},
$\Delta=tL-C\succeq0$. Decompose
\[
 L-L_{\infty}
 =\frac{\Delta}{t}
  -\frac{t-t_*}{tt_*}C.
\]
Both matrices on the right are positive semidefinite and have trace
$a=b\eta/(1+\eta)$. The triangle inequality for the nuclear norm
therefore gives $2a$. Each positive semidefinite term is bounded above
by $aI$, so their difference lies between $-aI$ and $aI$, proving the
operator norm estimate. Finally $\eta/(1+\eta)\le
\varepsilon/(1+\varepsilon)$.

For the second assertion expand the squared norm
\[
 \left\|L_H^{-1/2}C_H^{1/2}-\frac{S}{b}L_H^{1/2}\right\|_F^2
 =E_1(L)-2\frac{S^2}{b}+\frac{S^2}{b^2}\operatorname{tr}L_H
 =E_1(L)-\frac{S^2}{b}.
\]
Since $(L_1)_H=(b/S)C_H^{1/2}$, this is
\eqref{net:eq:residual}. No simultaneous diagonalization of $L$ and
$C$ is used. The Schatten Cauchy--Schwarz inequality
$\|XY\|_{S_1}\le\|X\|_F\|Y\|_F$ now gives
\begin{align*}
 \|L-L_1\|_{S_1}
 &\le\|L_H^{1/2}\|_F\,
      \|L_H^{-1/2}(L_H-(L_1)_H)\|_F\\
 &=\frac{b\sqrt b}{S}
       \sqrt{E_1(L)-S^2/b}
 \le b\sqrt\varepsilon.
\end{align*}
For clarity, the cited Schatten inequality follows from the dual formula
\[
 \|Z\|_{S_1}=\max_{U\text{ orthogonal}}
                 \operatorname{tr}(U^{\mathsf T}Z)
\]
and the ordinary Hilbert--Schmidt Cauchy--Schwarz inequality.
\end{proof}

These results settle the article's optimal-coupling question for
unrestricted graph support and positive input degrees. A prescribed sparse
support, individual edge caps, integral conductances, or a fixed input root
configuration introduces additional constraints. The universal lower
bounds in Theorem~\ref{net:thm:optima} remain valid under such restrictions,
but equality then requires feasibility of the unique displayed optimizer.
In particular, the first objective is a worst-case guarantee over roots;
the largest displacement for one fixed family can be smaller.

